\documentclass[10pt,a4paper]{amsart}
\usepackage{amsmath,amssymb,amsthm,graphicx}
\usepackage{verbatim}
\newenvironment{reference}{\comment}{\endcomment}
%\newenvironment{reference}{}{}
\newenvironment{slogan}{\comment}{\endcomment}
\newenvironment{history}{\comment}{\endcomment}

% For commutative diagrams we use Xy-pic
\usepackage[all]{xy}

% We use 2cell for 2-commutative diagrams.
\xyoption{2cell}
% Theorem environments.
%
\theoremstyle{plain}
\newtheorem{theorem}[subsection]{Theorem}
\newtheorem{proposition}[subsection]{Proposition}
\newtheorem{lemma}[subsection]{Lemma}

\theoremstyle{definition}
\newtheorem{definition}[subsection]{Definition}
\newtheorem{example}[subsection]{Example}
\newtheorem{exercise}[subsection]{Exercise}
\newtheorem{situation}[subsection]{Situation}

\theoremstyle{remark}
\newtheorem{remark}[subsection]{Remark}
\newtheorem{remarks}[subsection]{Remarks}

\numberwithin{equation}{subsection}

% Macros
%
\def\lim{\mathop{\mathrm{lim}}\nolimits}
\def\colim{\mathop{\mathrm{colim}}\nolimits}
\def\Spec{\mathop{\mathrm{Spec}}}
\def\Hom{\mathop{\mathrm{Hom}}\nolimits}
\def\Ext{\mathop{\mathrm{Ext}}\nolimits}
\def\SheafHom{\mathop{\mathcal{H}\!\mathit{om}}\nolimits}
\def\SheafExt{\mathop{\mathcal{E}\!\mathit{xt}}\nolimits}
\def\Sch{\mathit{Sch}}
\def\Mor{\mathop{\mathrm{Mor}}\nolimits}
\def\Ob{\mathop{\mathrm{Ob}}\nolimits}
\def\Sh{\mathop{\mathit{Sh}}\nolimits}
\def\NL{\mathop{N\!L}\nolimits}
\def\CH{\mathop{\mathrm{CH}}\nolimits}
\def\proetale{{pro\text{-}\acute{e}tale}}
\def\etale{{\acute{e}tale}}
\def\QCoh{\mathit{QCoh}}
\def\Ker{\mathop{\mathrm{Ker}}}
\def\Im{\mathop{\mathrm{Im}}}
\def\Coker{\mathop{\mathrm{Coker}}}
\def\Coim{\mathop{\mathrm{Coim}}}

% Boxtimes
%
\DeclareMathSymbol{\boxtimes}{\mathbin}{AMSa}{"02}

%
% Macros for moduli stacks/spaces
%
\def\QCohstack{\mathcal{QC}\!\mathit{oh}}
\def\Cohstack{\mathcal{C}\!\mathit{oh}}
\def\Spacesstack{\mathcal{S}\!\mathit{paces}}
\def\Quotfunctor{\mathrm{Quot}}
\def\Hilbfunctor{\mathrm{Hilb}}
\def\Curvesstack{\mathcal{C}\!\mathit{urves}}
\def\Polarizedstack{\mathcal{P}\!\mathit{olarized}}
\def\Complexesstack{\mathcal{C}\!\mathit{omplexes}}
% \Pic is the operator that assigns to X its picard group, usage \Pic(X)
% \Picardstack_{X/B} denotes the Picard stack of X over B
% \Picardfunctor_{X/B} denotes the Picard functor of X over B
\def\Pic{\mathop{\mathrm{Pic}}\nolimits}
\def\Picardstack{\mathcal{P}\!\mathit{ic}}
\def\Picardfunctor{\mathrm{Pic}}
\def\Deformationcategory{\mathcal{D}\!\mathit{ef}}
\usepackage{cleveref}
\crefname{originallemma}{original lemma}{original lemmas}
\Crefname{originallemma}{Original lemma}{Original lemmas}
\makeatletter
\newlabel{formal-spaces-lemma-completed-tensor-product}{{0GB4}{source}{}{sourcepointer-1}{}}
\newlabel{formal-spaces-lemma-completed-tensor-product@cref}{{[originallemma][0][]0GB4}{[1][0][]source}{}{}{}}
\makeatother
\crefname{originalreference}{original reference}{original references}
\Crefname{originalreference}{Original reference}{Original references}
\makeatletter
\newlabel{formal-spaces-situation-base-change-local-property}{{0GBC}{source}{}{sourcepointer-2}{}}
\newlabel{formal-spaces-situation-base-change-local-property@cref}{{[originalreference][0][]0GBC}{[1][0][]source}{}{}{}}
\makeatother
\crefname{originalreference}{original reference}{original references}
\Crefname{originalreference}{Original reference}{Original references}
\makeatletter
\newlabel{lemma-finite-type-red}{{0GC2}{source}{}{sourcepointer-3}{}}
\newlabel{lemma-finite-type-red@cref}{{[originalreference][0][]0GC2}{[1][0][]source}{}{}{}}
\makeatother
\title{Finite type after reduction is stable under base change}
\author{The Stacks Project Authors}
\date{stacks-project:restricted.tex:local-snapshot-20261012; GNU FDL 1.2 or later}
\begin{document}
\maketitle
\noindent\textit{Editorial note.} All delivered mathematical text is the author's own, copied from the pinned Stacks Project chapter file restricted.tex; no statement or proof step was rewritten, reordered or removed.\par\medskip
\noindent\textit{Editorial note.} The counted proof is the author's complete argument: the proof identifies the reduction of the completed tensor product as the ordinary base change of the reductions and applies stability of finite type under base change\par\medskip
\noindent\textit{Editorial note.} The label lemma-base-change-finite-type-red names the counted statement, and the proof environment immediately following it in the source is delivered verbatim as its proof.\par\medskip
\noindent\textit{Editorial note.} The counted statement is carried by the author's own lemma environment in the chapter source.\par\medskip
\noindent\textit{Allowed standard prerequisites (editorial).}\begin{itemize}
\item the property P on arrows of WAdm\textasciicircum{}count defined in lemma-finite-type-red of the same chapter: B/b -> A/a is of finite type for the ideals of topologically nilpotent elements
\item weak admissibility makes the ideals of topologically nilpotent elements open
\item the completed tensor product of weakly admissible rings and its kernel being a weak ideal of definition (compare Formal Spaces, Lemma formal-spaces-lemma-completed-tensor-product)
\item finite type of ring maps is stable under base change
\end{itemize}\medskip
\ifdefined\hypertarget\hypertarget{sourcepointer-1}{}\fi
\noindent\textit{Source reference.} Original lemma 0GB4: sources/stacks-chow/tags.txt, source lines 20359--20359. Prerequisite pointer invoked at proof line(s) 2719; it names a standard fact of a sibling chapter that the counted argument uses, and every step of the counted proof built on it is retained verbatim in the delivered body.\par\medskip
\ifdefined\hypertarget\hypertarget{sourcepointer-2}{}\fi
\noindent\textit{Source reference.} Original reference 0GBC: sources/stacks-chow/tags.txt, source lines 20367--20367. Prerequisite pointer invoked at proof line(s) 2701; it names a standard fact of a sibling chapter that the counted argument uses, and every step of the counted proof built on it is retained verbatim in the delivered body.\par\medskip
\ifdefined\hypertarget\hypertarget{sourcepointer-3}{}\fi
\noindent\textit{Source reference.} Original reference 0GC2: sources/stacks-restricted/pinned-original-source.tex, source lines 2631--2631. Prerequisite pointer invoked at proof line(s) 2699; it names a standard fact of the same chapter that the counted argument uses, and every step of the counted proof built on it is retained verbatim in the delivered body.\par\medskip
\begin{lemma}
\label{lemma-base-change-finite-type-red}
Consider the property $P$ on arrows of $\textit{WAdm}^{count}$ defined in
Lemma \ref{lemma-finite-type-red}. Then $P$ is stable under base change
(Formal Spaces, Situation
\ref{formal-spaces-situation-base-change-local-property}).
\end{lemma}

\begin{proof}
The statement makes sense by Lemma \ref{lemma-finite-type-red}.
To see that it is true assume we have morphisms $B \to A$ and $B \to C$
in $\textit{WAdm}^{count}$ such that $B/\mathfrak b \to A/\mathfrak a$
is of finite type where $\mathfrak b \subset B$ and $\mathfrak a \subset A$
are the ideals of topologically nilpotent elements.
Since $A$ and $B$ are weakly admissible, the ideals
$\mathfrak a$ and $\mathfrak b$ are open.
Let $\mathfrak c \subset C$ be the (open) ideal
of topologically nilpotent elements. Then we find a surjection
$A \widehat{\otimes}_B C \to
A/\mathfrak a \otimes_{B/\mathfrak b} C/\mathfrak c$
whose kernel is a weak ideal of definition and hence consists
of topologically nilpotent elements
(please compare with the proof of Formal Spaces,
Lemma \ref{formal-spaces-lemma-completed-tensor-product}). Since already
$C/\mathfrak c \to A/\mathfrak a \otimes_{B/\mathfrak b} C/\mathfrak c$
is of finite type as a base change of $B/\mathfrak b \to A/\mathfrak a$
we conclude.
\end{proof}

\end{document}
